\documentclass[11pt]{article}
\usepackage[margin=1in]{geometry}
\usepackage{amsmath,amssymb,amsthm,booktabs,array}
\usepackage[T1]{fontenc}
\usepackage{lmodern}
\usepackage[colorlinks=true,linkcolor=blue,citecolor=blue,urlcolor=blue]{hyperref}
\newtheorem{theorem}{Theorem}
\newtheorem{lemma}[theorem]{Lemma}
\newtheorem{corollary}[theorem]{Corollary}
\newcommand{\R}{\mathbb R}
\newcommand{\C}{\mathbb C}
\newcommand{\N}{\mathbb N}
\newcommand{\ph}{\varphi}
\title{Global convexity of the Bernstein basis square-sum\\
\large A finite Fourier proof and a Lean formalization}
\author{Myeongjin Kang}
\date{25 October 2026}
\setlength{\emergencystretch}{2em}
\begin{document}
\maketitle
\begin{abstract}
For every nonnegative integer $n$, the polynomial
$\ph_n(x)=\sum_{k=0}^n\binom nk^2x^{2k}(1-x)^{2(n-k)}$
is convex on the entire real line. We give a finite Fourier proof: its defining
sum is exactly an average of the convex functions
$x\mapsto |1-x+x\zeta|^{2n}$ over the $(n+1)$st roots of unity.
The accompanying Lean module proves the averaging identity, global convexity,
symmetry, monotonicity on either side of $1/2$, and the sharp midpoint minimum.
This formalizes a known result; no mathematical novelty or priority is claimed.
The related MathDB entries 327180--327182 constitute one outcome.
\end{abstract}

\section{Statement, sources, and scope}
Put
\[
 b_{nk}(x)=\binom nk x^k(1-x)^{n-k},\qquad
 \ph_n(x)=\sum_{k=0}^n b_{nk}(x)^2 .
\]
Here $n\in\N$ includes zero and $x$ is real. Powers with exponent zero
have their usual polynomial meaning, including at the endpoints.

Gonska, Ra\c{s}a, and Rusu state convexity on $[0,1]$ as Conjecture 3.2,
the two monotonicity assertions as Conjecture 3.3, and the midpoint minimum
as Conjecture 3.4 in \cite{GRR}. The retrieved arXiv v1 itself subsequently
proves the minimum in Lemma 3.5, crediting Thorsten Neuschel. Consequently,
the frozen MathDB 327182 description that the source supplies no resolution
is inaccurate. We preserve the source wording in the provenance record and note this correction.

Nikolov \cite{Nikolov} proves the convexity conjecture in Corollary 4,
through inequalities for ultraspherical polynomials. Gavrea and Ivan
\cite{GI} give a separate representation-based proof in 2014; their
published abstract and section descriptions explicitly identify Parseval's
identity as a tool. The full text of that article was not obtained for this
artifact, so we do not attribute a particular numbered identity to it.
Our finite Fourier proof is in that established methodological family.
We claim neither a new theorem nor the first formalization.

\begin{theorem}\label{thm:main}
For every $n\in\N$, $\ph_n$ is convex on $\R$. In particular it is
convex on $[0,1]$, as requested in MathDB 327180.
\end{theorem}
The global statement concerns the same polynomial extension, not a
probability interpretation outside $[0,1]$.

\section{The exact finite averaging identity}
The next argument provides the bridge from the original square-sum;
no alternate representation is assumed as a hypothesis.

\begin{lemma}[Finite Parseval identity]\label{lem:parseval}
Let $N\geq1$, let $a_0,\ldots,a_{N-1}$ be real numbers, and let
$\zeta_j=\exp(2\pi i j/N)$. Then
\[
 \sum_{j=0}^{N-1}\left|\sum_{k=0}^{N-1}a_k\zeta_j^k\right|^2
 =N\sum_{k=0}^{N-1}a_k^2.
\]
\end{lemma}
\begin{proof}
For $0\leq k,l<N$, the finite geometric sum gives
\[
 \sum_{j=0}^{N-1}\zeta_j^k\overline{\zeta_j^l}
 =\sum_{j=0}^{N-1}\exp(2\pi i j(k-l)/N)
 =\begin{cases}N,&k=l,\\0,&k\ne l.\end{cases}
\]
Indeed, when $k\ne l$ the ratio is not one, while its $N$th power is
one. Expand the squared absolute value as the product with its complex
conjugate and interchange the three finite sums. The off-diagonal terms
vanish, leaving $N\sum_k a_k^2$.
\end{proof}

\begin{lemma}[Averaging the Bernstein square-sum]\label{lem:average}
For every $n\in\N$ and $x\in\R$, with $N=n+1$,
\begin{equation}\label{eq:average}
 \ph_n(x)=\frac1{n+1}\sum_{j=0}^{n}
       |1-x+x\zeta_j|^{2n}.
\end{equation}
\end{lemma}
\begin{proof}
The binomial theorem gives, for each $j$,
\[
 (1-x+x\zeta_j)^n
 =\sum_{k=0}^n\binom nk x^k(1-x)^{n-k}\zeta_j^k
 =\sum_{k=0}^n b_{nk}(x)\zeta_j^k.
\]
Apply Lemma~\ref{lem:parseval} to $a_k=b_{nk}(x)$, use
$|z^n|^2=|z|^{2n}$, and divide by $n+1>0$.
For $n=0$, both sides of \eqref{eq:average} equal one, so no exceptional
case or division by $n$ is present.
\end{proof}

\section{Convexity and the shape corollaries}
\begin{proof}[Proof of Theorem~\ref{thm:main}]
For any fixed $z\in\C$, the function $x\mapsto |1-x+xz|$ is convex
on $\R$, being the norm composed with a real affine map. Its values
are nonnegative. For every nonnegative integer $m$, the function
$u\mapsto u^m$ is convex and nondecreasing on $[0,\infty)$; the case
$m=0$ is constant. Thus $x\mapsto |1-x+xz|^{2n}$ is convex.
Equation~\eqref{eq:average} expresses $\ph_n$ as a finite positive
average of such functions, proving the result for every real $x$.
\end{proof}

\begin{corollary}\label{cor:shape}
For every $n\in\N$, the following hold:
\begin{enumerate}
\item $\ph_n(1-x)=\ph_n(x)$ for every $x\in\R$;
\item $\ph_n$ is nonincreasing on $(-\infty,1/2]$ and nondecreasing
on $[1/2,\infty)$;
\item for every $x\in\R$,
\begin{equation}\label{eq:minimum}
 \ph_n(x)\geq\ph_n(1/2)=4^{-n}\binom{2n}{n}.
\end{equation}
\end{enumerate}
In particular these imply the assertions of MathDB 327181 and 327182
on their stated intervals. Non-strict monotonicity accommodates $n=0$.
\end{corollary}
\begin{proof}
The substitution $k\mapsto n-k$ in the defining sum proves symmetry.
Convexity and symmetry give
\[
 \ph_n(1/2)=\ph_n\bigl((x+(1-x))/2\bigr)
 \leq\tfrac12\ph_n(x)+\tfrac12\ph_n(1-x)=\ph_n(x).
\]
If $x<y<1/2$, express $y$ as a strict convex combination of $x$ and
$1/2$. The resulting convexity inequality, together with
$\ph_n(1/2)\leq\ph_n(y)$, gives $\ph_n(y)\leq\ph_n(x)$.
The cases $x=y$ and $y=1/2$ follow directly; symmetry gives monotonicity
on the right half-line.

At the midpoint, $b_{nk}(1/2)=2^{-n}\binom nk$. Vandermonde's identity
$\sum_{k=0}^n\binom nk^2=\binom{2n}{n}$ then proves the equality
in \eqref{eq:minimum}.
\end{proof}

\section{Independent verification of the proposed coefficient identity}
The proposed positive even-power expansion is correct:
\begin{equation}\label{eq:positive}
 \boxed{\displaystyle
 \ph_n(x)=4^{-n}\sum_{j=0}^n
 \binom{2j}{j}\binom{2n-2j}{n-j}(2x-1)^{2j}.}
\end{equation}
Here is a derivation independent of the finite Fourier implementation.
Continuous Fourier orthogonality, proved by integrating the finite
exponential sums term by term, gives
\[
 \ph_n(x)=\frac1{2\pi}\int_0^{2\pi}|1-x+xe^{i\theta}|^{2n}\,d\theta.
\]
Set $t=2x-1$. Direct multiplication shows
\[
 |1-x+xe^{i\theta}|^2
 =\cos^2(\theta/2)+t^2\sin^2(\theta/2).
\]
After $\theta=2u$ and reflection about $\pi/2$, this becomes
\begin{equation}\label{eq:betaexpansion}
 \ph_n(x)=\sum_{j=0}^n\binom nj t^{2j}J_{j,n-j},\qquad
 J_{a,b}=\frac2\pi\int_0^{\pi/2}\sin^{2a}u\cos^{2b}u\,du.
\end{equation}
For nonnegative integers $a,b$, elementary integration by parts gives
\[
 J_{a,b}=\frac{(2a)!(2b)!}{4^{a+b}a!b!(a+b)!}.
\]
For completeness, $J_{0,0}=1$, symmetry gives $J_{a,b}=J_{b,a}$,
and for $b\geq1$ the integral of the derivative of
$\sin^{2a+1}u\cos^{2b-1}u$ has zero boundary terms. Replacing
$\sin^2u$ by $1-\cos^2u$ yields
$2(a+b)J_{a,b}=(2b-1)J_{a,b-1}$.
These relations establish the displayed factorial expression by induction.
Multiplying it by $\binom nj$ in \eqref{eq:betaexpansion} proves
\eqref{eq:positive} for all $n$, including zero.

For example, $n=2$ gives
\[
 \ph_2(x)=\frac38+\frac14t^2+\frac38t^4.
\]
All coefficients in \eqref{eq:positive} are nonnegative, furnishing another
proof of global convexity. This appendix is a mathematical verification;
the closed coefficient formula and its integrals are not declarations in
the Lean module. They are not assumptions or missing bridges in the
formal convexity proof, which instead proves \eqref{eq:average} exactly.
An additional integer-coefficient regression checks \eqref{eq:positive}
for $0\leq n\leq32$; that finite computation is not used as a proof for all $n$.

\section{Formalization and verification contract}
The single self-contained module is \nolinkurl{formal/BernsteinShape.lean}.
It imports the pinned mathlib and has no dependencies on other campaign
outcomes. In the formal proof, characters on $\mathbb Z/(n+1)\mathbb Z$
replace explicit complex exponentials. The orthogonality lemma uses
mathlib's primitive standard character and its character-sum theorem.
The coefficient indices are $\operatorname{Fin}(n+1)$; the proof that
their residues are distinct uses their strict bounds. This explicitly
prevents frequency aliasing. The binomial theorem then supplies the
original coefficients, including their signs when $x\notin[0,1]$.

\begin{center}
\small
\begin{tabular}{@{}p{0.44\textwidth}p{0.49\textwidth}@{}}
\toprule
Lean declaration & Content \\
\midrule
\nolinkurl{character_orthogonality} & Exact cancellation for unequal indices \\
\nolinkurl{finite_parseval} & Real coefficient square-sum identity \\
\nolinkurl{binomial_fourier} & Original Bernstein coefficients in the Fourier sum \\
\nolinkurl{finite_average_identity} & Equation \eqref{eq:average}, for all real $x$ \\
\nolinkurl{convexOn_univ} & Global convexity for every $n\in\N$ \\
\nolinkurl{bernstein_square_sum_convex} & Exact expanded MathDB 327180 formula on $[0,1]$ \\
\nolinkurl{symmetry}, \nolinkurl{midpoint_minimum} & Symmetry and global minimum \\
\nolinkurl{antitoneOn_left}, \nolinkurl{monotoneOn_right} & Half-line monotonicity \\
\nolinkurl{midpoint_value}, \nolinkurl{sharp_lower_bound} & Exact sharp constant \\
\bottomrule
\end{tabular}
\end{center}

The compiler is Lean 4.34.1. The pinned mathlib revision is
\begin{center}\small\ttfamily
d13f23b723b8a846827a245b89c10fc7d3f11612
\end{center}
The remaining dependency revisions are recorded in \nolinkurl{pins.json}.
The campaign verifier compiles
a byte-identical isolated source copy with warnings treated as errors,
checks every declared definition and theorem, and reports its axioms.
Only \nolinkurl{propext}, \nolinkurl{Classical.choice}, and \nolinkurl{Quot.sound}
are permitted. The verification receipts, source hashes, and complete
compiler output accompany this artifact. Independent review accompanies the release and remains separate from successful kernel verification. The completed audit checks
20 declarations. The release includes the campaign verifier and pinned dependency revisions; its fresh-source compilation and axiom audit are reproducible.

\begin{thebibliography}{9}
\bibitem{GRR}
H. Gonska, I. Ra\c{s}a, and M.-D. Rusu,
\emph{Chebyshev--Gr\"uss-type inequalities via discrete oscillations},
arXiv:1401.7908v1 (2014), Conjectures 3.2--3.4 and Lemma 3.5.
\url{https://arxiv.org/abs/1401.7908v1}.

\bibitem{Nikolov}
G. Nikolov,
\emph{Inequalities for ultraspherical polynomials. Proof of a conjecture
of I. Ra\c{s}a}, arXiv:1402.6539v1 (2014), Corollary 4;
Journal of Mathematical Analysis and Applications \textbf{418} (2014),
852--860. \url{https://doi.org/10.1016/j.jmaa.2014.04.022}.

\bibitem{GI}
I. Gavrea and M. Ivan,
\emph{On a conjecture concerning the sum of the squared Bernstein
polynomials}, Applied Mathematics and Computation \textbf{241} (2014),
70--74. \url{https://doi.org/10.1016/j.amc.2014.05.003}.
\end{thebibliography}
\end{document}
